\section{Scope and System Context}
\label{sec:scope}

Table~\ref{tab:scope} declares what this project will and will
not do, in pairs; Fig.~\ref{fig:context} shows the system as a
single block with the actors it exchanges information with.

\begin{table}[h]
\centering
\caption{Project scope, declared in pairs.}
\label{tab:scope}
\begin{tabular}{p{3.7cm}p{3.7cm}}
\toprule
\textbf{In scope} & \textbf{Out of scope} \\
\midrule
Sense motion at the wrist (watch IMU) and at the chest and/or
thigh (tethered IMUs on one custom board) &
Other body positions, and sensors other than IMUs \\
Fuse wrist, chest and/or thigh features to classify sedentary,
walking or running, falling back to fewer sensors on dropout &
Fall detection, and activity types beyond these three
(including vehicles) \\
Report each activity change to a companion computer over Wi-Fi,
with retries after an outage &
Bluetooth or cellular connectivity \\
Computer-side logging and a simple progress view &
Long-term cloud storage and multi-user accounts \\
On-watch display of basic status &
Enclosure design and battery-life optimisation \\
Low-fidelity prototype: off-the-shelf modules on breadboard or
perfboard &
Fabricating the custom PCB (designed and built in Project II) \\
\bottomrule
\end{tabular}
\end{table}

\begin{figure}[h]
\centering
\begin{tikzpicture}[
  font=\footnotesize,
  sysb/.style={draw, thick, rectangle, minimum width=2.6cm, minimum height=1.1cm, align=center},
  actor/.style={draw, ellipse, dashed, minimum width=2.1cm, minimum height=0.8cm, align=center, inner sep=1pt},
  stake/.style={draw, ellipse, dotted, thick, minimum width=2.3cm, minimum height=0.8cm, align=center, inner sep=1pt, font=\scriptsize},
  lbl/.style={font=\scriptsize, align=center},
  >=stealth
]
\node[sysb] (sys) {Activity\\Monitoring System};
\node[actor, above=1.3cm of sys] (wearer) {Wearer};
\node[actor, right=2.0cm of sys] (wifi) {Wi-Fi\\Network};
\node[actor, below=1.1cm of wifi] (pc) {Companion\\computer};
\node[stake, below=1.1cm of sys] (family) {Family member /\\support person};
\node[stake, below=0.8cm of pc] (employer) {Employer or insurer\\(wellness program)};

\draw[->] ([xshift=-0.55cm]sys.north |- wearer.south) -- node[left, lbl] {motion (wrist,\\chest, thigh)} ([xshift=-0.55cm]sys.north);
\draw[->] ([xshift=0.55cm]sys.north) -- node[right, lbl] {activity\\feedback} ([xshift=0.55cm]sys.north |- wearer.south);
\draw[<->] (sys) -- node[above, lbl] {activity reports /} node[below, lbl] {acknowledgements} (wifi);
\draw[<->] (wifi) -- node[right, lbl] {degraded or\\cut in testing} (pc);
\draw[->, dotted, thick] (pc) -- node[above, lbl] {weekly trends} (family);
\draw[->, dotted, thick] (pc) -- node[right, lbl] {aggregate \%\\only} (employer);
\end{tikzpicture}
\caption{Context diagram. The watch and custom hardware sit inside
the single system boundary; the dashed ellipses are the actors
behind each arrow, and every solid arrow crossing the boundary is
something the project must build and defend. Only classified
activity reports leave the system; raw sensor data does not.
Dotted ellipses are the other stakeholders from the Project Idea
Form; their views are not built in this PoC.}
\label{fig:context}
\end{figure}

\textbf{Stakeholder collision:} the wearer wants feedback that
feels current, with minimal personal effort (no manual logging,
low battery drain), while the shared Wi-Fi network favours low,
predictable traffic. Streaming sensor data or sending frequent
updates would make the progress view feel more live, but would
directly compete with keeping network load light. This collision
is resolved in Section~\ref{sec:reqs} by sending a report only
when the activity changes, within a stated time limit, rather than
streaming continuously.
